\section{Case~2: When 75.56\% accuracy was worse than a one-class baseline}
\label{app:case-heart-disease}

\subsection{Why VA became the review question}
The case began with a familiar multi-cohort task: establish comparable local classification baselines for the Cleveland, Hungarian and VA partitions of the UCI Heart Disease dataset, then review the results together. The partitions contain related clinical variables but differ in sample size, missingness and outcome prevalence. The working question was therefore not which row in a score table was ``best'', but whether a common analysis protocol and cross-review would reveal a result whose interpretation changed with cohort composition.

One operator enacted three roles through separate SciForge accounts on one Windows device. Role~B analysed Cleveland and maintained the shared research record. Role~C analysed Hungarian and reviewed the submissions. Role~D analysed VA. Each role worked in a separate directory. Within the project, the roles exchanged method descriptions, aggregate outputs and review comments rather than patient-level rows. This arrangement made the handoff and interpretation reviewable while keeping the three analyses distinct.

\subsection{From three local fits to one reviewable record}
The roles used one custom NumPy baseline. The script converted the multiclass disease field to a binary target, removed rows that could not be processed under the shared schema and applied a fixed categorical vocabulary. Each partition used random seed 43, a class-stratified train/test split and an $L_2$-regularised logistic-regression model fitted on CPU. Continuous-feature means and standard deviations were estimated from the training split only. Every local report recorded the retained sample count, train/test sizes, retained-cohort positivity, test AUROC and test accuracy.

Roles~C and D submitted the Hungarian and VA summaries to Role~B through project tasks. Role~B checked the declared protocol, accepted both materials and placed the three baselines in the shared record. The task board and the shared note formed the operation record shown in Fig.~\ref{fig:case-heart-operations}: the first view documents the handoff of local reports, and the second documents the discussion that changed how the VA row was read.

\begin{figure}[htbp]
  \centering
  \begin{minipage}{\linewidth}
    \centering
    \includegraphics[width=\linewidth,trim=390 450 220 0,clip]{cases/case-02-heart-disease/screenshots/project-handoff.png}
    \par\smallskip\noindent\textbf{a} Project handoff and accepted local reports.
  \end{minipage}

  \vspace{0.8em}
  \begin{minipage}{\linewidth}
    \centering
    \includegraphics[width=0.75\linewidth,trim=380 0 0 0,clip]{cases/case-02-heart-disease/screenshots/metric-review.png}
    \par\smallskip\noindent\textbf{b} Shared-note review of the VA metrics.
  \end{minipage}
  \caption{SciForge records of the local-report handoff and the subsequent review of the VA result.}
  \label{fig:case-heart-operations}
\end{figure}

\subsection{How the original review changed the report}
During the original 12--13 November 2026 demonstration, Cleveland and Hungarian returned AUROC values above 0.87. VA looked less straightforward: its test accuracy was 0.7556, yet its AUROC was only 0.5629. Role~C asked whether the apparently usable accuracy mainly reflected the dominant class rather than reliable separation of positive and negative cases. The question changed the report in three concrete ways. Accuracy and AUROC were displayed together; the positivity column was explicitly labelled as a statistic over all retained samples rather than the test set; and repeated splits or confidence intervals were named as the next analysis. The collaborators did not use the three test sets to rank the models because each partition had a different composition.

The original demonstration did not retain the VA test-label counts or thresholded predictions. It therefore supported the discrepancy---0.7556 accuracy versus 0.5629 AUROC---and the decision to investigate prevalence, but not a numerical explanation of that discrepancy. The screenshots in Fig.~\ref{fig:case-heart-operations} document this original stage. The detailed counts below were reconstructed later and were not written back into the original project record.

\subsection{What the 2026-11-22 reconstruction added}
The original aggregate results are summarised in Table~\ref{tab:case-heart-baselines}. Cleaning retained all 303 Cleveland rows, 261 of 294 Hungarian rows and 130 of 200 VA rows. Positivity was calculated over all retained samples after cleaning, whereas AUROC and accuracy were calculated on each partition's test split.

\begin{table}[htbp]
  \centering
  \caption{Local logistic-regression baselines across the three UCI Heart Disease partitions. Positivity refers to all retained samples; AUROC and accuracy refer to the local test set.}
  \label{tab:case-heart-baselines}
  \begin{tabular}{lrrrrr}
    \toprule
    Partition & Retained & Train / test & Positivity (\%) & Test AUROC & Test accuracy \\
    \midrule
    Cleveland & 303 & 199 / 104 & 45.87 & 0.8787 & 0.7788 \\
    Hungarian & 261 & 171 / 90 & 37.55 & 0.9044 & 0.8333 \\
    VA        & 130 & 85 / 45   & 77.69 & 0.5629 & 0.7556 \\
    \bottomrule
  \end{tabular}
\end{table}

On 2026-11-22, the hash-verified script was rerun against the public source data to reconstruct the missing VA diagnostic. Of 45 test cases, 35 were positive (77.78\%) and 10 were negative. At a threshold of 0.5, the model predicted 38 positives (84.44\%) and seven negatives. It produced 31 true positives, four false negatives, three true negatives and seven false positives, for 34 correct classifications. Thus, the model's 75.56\% accuracy was below the 77.78\% obtained by assigning every test case to the positive class. AUROC remained 0.5629. Figure~\ref{fig:case-heart-disease} keeps the two evidence stages visible: panel~a shows the original collaborative path, while panel~b shows the later diagnostic that tested the review question.

\begin{figure}[htbp]
  \centering
  \includegraphics[width=\textwidth]{cases/case-02-heart-disease/figures/fig5_heart_disease.pdf}
  \caption{Cross-review and follow-up analysis of the VA result. \textbf{a}, Original SciForge workflow from local analyses to cross-review and revision. \textbf{b}, Subsequent VA diagnostic reconstructed from the versioned analysis.}
  \label{fig:case-heart-disease}
\end{figure}

\subsection{What the measurements support}
The VA result supports a narrow conclusion: its test accuracy was not evidence that the model improved on a trivial majority-class decision, and its AUROC indicated weak ranking performance on this split. The confusion counts also show limited recognition of negative cases, with three true negatives and seven false positives. These observations explain why accuracy had to be interpreted beside class composition and AUROC. They do not identify the mechanism behind the weak discrimination. Missing-data handling, cohort composition, feature shift, model specification and split variability remain competing explanations.

That boundary leads directly to the next request rather than to a stronger claim. A follow-up experiment should harmonise preprocessing across partitions, repeat the stratified split over multiple seeds, report uncertainty for AUROC and accuracy, and compare threshold-dependent errors before any local-versus-collaborative training claim. In this sense, cross-review produced a useful research transition: it turned one plausible-looking number into a falsifiable question and a defined measurement plan.

\subsection{Scope, provenance and availability}
The evidence chain is reproducible but deliberately limited. The demonstration used one operator, one Windows device and three test accounts; it did not validate physical institutional separation or a cryptographic data boundary. The model was a custom NumPy implementation, not the official FLamby benchmark. No federated training, joint-versus-local comparison, repeated-split analysis, confidence interval or external clinical validation was performed. None of the outputs is intended for clinical decision-making.

The three source partitions are available from the UCI Machine Learning Repository, Heart Disease dataset, DOI: \href{https://doi.org/10.24432/C52P4X}{10.24432/C52P4X}. The analysis script, original aggregate results, SHA-256 provenance, supplementary VA audit and reproduction command are available in the \href{https://github.com/MultiHugo/sciforge-heart-disease-collaboration}{public Heart Disease case repository}. The repository reproduces the numerical evidence without SciForge credentials.
